\section{由收缩循环得到的局部识别}
\label{sec:shrinking}

令 $z=(x,a)$，并固定 $z_0=(x_0,a_0)$。在局部设计 $n$ 中，令
\[
z_{1n},\ldots,z_{M_n n}
\]
为抽样设计点，并写
\[
A(z):=(V\circ\Gamma_A)(z),
\qquad
B(z):=(V\circ\Gamma_B)(z),
\qquad
H(z):=A(z)-B(z).
\]
把 $Q_n$ 中的延续标签按如下次序排列：
\[
A(z_{1n}),\ldots,A(z_{M_n n}),
B(z_{1n}),\ldots,B(z_{M_n n}).
\]
当有限表示通过结构模型解释时，延续势 $v$（在共同归一化下）把这些标签分别赋值为 $A(z_{in})$ 与 $B(z_{in})$。为了给出锐利的局部结论，我们直接使用满足观测循环方程的约化式函数 $A$ 与 $B$。设计 $n$ 中的原始补偿比较具有观测当前关联 $D_n$、延续对比 $Q_n$ 与精确转移 $\tau_n$。

\subsection{目标环境与观测比较}
\label{sec:environment-support}

目标环境在 $z_0$ 的一个邻域上定义为
\[
W_B(x,a)=u_B(x,a)+V(\Gamma_B(x,a)).
\]
令 $\mathcal E_n$ 表示局部设计 $n$ 中观测到的补偿比较。族 $\{\mathcal E_n\}$ 不必包含由环境 $B$ 中所有邻近方案所组成的比较。它还可以包含这样的诱导菜单：把一个当前后果 $\xi$ 与某个延续合同 $\Gamma_J(z)$ 配对，即使该复合方案并不属于目标选择集。在保险例子中，这类合同变体保持当前保障不变，却替换成不同的状态依赖续保条款。下文的识别均相对于 $\{\mathcal E_n\}$ 而言；$\Gamma_B$ 仍在完整的局部选择集上有定义。

\subsection{来自循环空间的延续对比}

令
\[
S_n:=(I_{M_n}\ I_{M_n}),
\qquad
T_n:=(I_{M_n}\ 0).
\]
与 $A$--$B$ 对比相关的整数循环格为
\begin{equation}
\mathcal C_n^{\mathbb Z}
:=
\left\{
r\in\mathbb Z^{E_n}:
D_nr=0,\ S_nQ_nr=0
\right\}.
\label{eq:integer-local-cycles}
\end{equation}
它们刻画观测比较之间的有限复制与消去关系。

这些观测方程可以线性组合。定义实循环空间
\begin{equation}
\mathcal C_n
:=
\operatorname{span}_{\mathbb R}\mathcal C_n^{\mathbb Z}
=
\left\{
r\in\mathbb R^{E_n}:
D_nr=0,\ S_nQ_nr=0
\right\}.
\label{eq:real-local-cycles}
\end{equation}
之所以有上述等式，是因为 $D_n$ 与 $Q_n$ 的系数都是整数：有理零空间向量张成实零空间，而通过清除分母可把它们转化为整数循环。

对 $r\in\mathcal C_n$，总延续对比具有形式
\[
Q_nr=(\lambda,-\lambda),
\qquad
\lambda=T_nQ_nr.
\]
由此得到的对比空间为
\begin{equation}
\boxed{
\Lambda_n:=T_nQ_n\mathcal C_n
\subseteq\mathbb R^{M_n}.}
\label{eq:gap-weight-space}
\end{equation}
在消去条件下，每个 $r\in\mathcal C_n^{\mathbb Z}$ 都满足
\[
r^{\mathsf T}\tau_n
=
\sum_i\lambda_iH(z_{in}).
\]
由线性性，同一恒等式对任意 $r\in\mathcal C_n$ 也成立：
\begin{equation}
r^{\mathsf T}\tau_n
=
\sum_i\lambda_iH(z_{in}),
\qquad
Q_nr=(\lambda,-\lambda).
\label{eq:gap-cycle-price}
\end{equation}
